\documentclass{article}

\usepackage{tabularx}
\usepackage{booktabs}

\title{Problem Statement and Goals\\\progname}

\author{\authname}

\date{}

\input{../Comments}
%% Common Parts

\newcommand{\progname}{Cut Edge Weight: Data Separability Measurement} % PUT YOUR PROGRAM NAME HERE
\newcommand{\authname}{Lesley Wheat} % AUTHOR NAMES                  

\usepackage{hyperref}
    \hypersetup{colorlinks=true, linkcolor=blue, citecolor=blue, filecolor=blue,
                urlcolor=blue, unicode=false}
    \urlstyle{same}
                                


\begin{document}

\maketitle

\begin{table}[hp]
\caption{Revision History} \label{TblRevisionHistory}
\begin{tabularx}{\textwidth}{llX}
\toprule
\textbf{Date} & \textbf{Developer(s)} & \textbf{Change}\\
\midrule
2023-01-20 & Lesley Wheat & Creation\\
\bottomrule
\end{tabularx}
\end{table}

\section{Problem Statement}

Classification is a common machine learning problem and there are many techniques for solving classification problems. 
Unfortunately, not all classification problems are easy and it can be difficult to explain why classifiers preform poorly on certain datasets. 
Data complexity has been defined as a characteristic of the data that relates to the difficulty of the classification problem. \cite{luengo_automatic_2015} 
To further refine this definition, we consider data complexity to be represented by aspects of the data that introduce difficulty for the classifiers. 

One of these aspects is separability, an intrinsic characteristic that describes how the data points from different classes mix together. \cite{guan_novel_2021} 
A classic measurement for separability is Fraction of Points on Class Boundary.
Based on the work of Friedman and Rafsky \cite{friedman_multivariate_1979} and documented by Basu and Ho\cite{ho_complexity_2002}, this measurement constructs a minimum spanning tree over the dataset then calculates the number of points with a connection to another class over all points.

The Cut Edge Weight measurement \cite{zighed_statistical_2005} is a variation of the Fraction of Points on Class Boundary, proposed by Zighed et al. 
The algorithm for the normalized measurement begins by constructing a Toussaint’s Relative Neighbourhood Graph \cite{toussaint_relative_1980} over the entire dataset then computing the number of edges which connect different classes over the total number of edges. 
This measurement is a viable way of measuring separability because it has a high coefficient of determination ($R^2$) with the mean error rate of several classifiers.
This algorithm does not currently have an available open source implementation, unlike others listed in \cite{ho_complexity_2002}.

\subsection{Problem}

The problem is the computation of the Cut Edge Weight measurement over a dataset.

\subsection{Inputs and Outputs}

\subsubsection{Inputs}

The program should receive the dataset as an input. 
The dataset should contain samples with features (in machine learning, a feature is an individual measurable property or characteristic \cite{bishop_pattern_2006}). 
Each sample should have the same number of features. 
Each feature should have a numerical value.

\subsubsection{Outputs}

The program should output the value of the normalized Cut Edge Weight measurement for the input dataset as a floating point number.

\subsection{Stakeholders}

The stakeholders are machine learning researchers.

\subsection{Environment}

The program should run on a Windows 11 operating system. 
The program should run on a laptop and should not require a GPU.

\section{Goals}

The goal of the software is to calculate the value of the normalized Cut Edge Weight Statistic as defined in \cite{zighed_statistical_2005} using a Relative Neighbourhood Graph with no weighting.


%\pagebreak

\section{Stretch Goals}

Stretch goals include:

\begin{itemize}
\item The option to weigh the edge calculation using the distance measurement of $(1+d_{ij})^{-1}$ where $d_{ij}$ is the Euclidean distance between points $i$ and $j$.
\item The option to weigh the edge calculation using rank so that the weight is $r_i^{-1}$. Where  $r_j$ is calculated as the rank of vertex $j$ among the neighbours of vertex $i$.
\item Generate the p-value of the statistic as defined in \cite{zighed_statistical_2005}.
\item Give the user the option to run the algorithm using a Gabrial Graph instead of Relative Neighbourhood Graph.
\item Give the user the option to run the algorithm using a Minimal Spanning Tree instead of Relative Neighbourhood Graph.
\item Give the user the option to output the statistic without normalization.
\item Package the code in a way accessible to the public.
\end{itemize}


\bibliographystyle{IEEEtran}
\bibliography{refs/Ref-ProbState}

\end{document}